\subsection{R-31: a scalar-period obstruction in the next-row two-product family}

\paragraph{Statement R-31.}
Let $m\geq3$, $a\geq1$, $n=m+a$, and $r=2m+1$. Define the
common two-product padded cone
\[
 \mathcal T_{m,n,r}
 =\overline{\left\{\ell^a\left(\prod_{i=1}^m L_i+
                         \prod_{i=1}^m M_i\right)\right\}}
                       \subseteq\operatorname{Sym}^n\mathbb C^r.
\]
The factors are arbitrary linear forms. This cone lies in both
projected permanent and determinant cones by the cycle and
padding constructions. Suppose $s,b$ are positive integers with
$ms=rb$, and put
\[
 \nu=(as+b,b^{,r-1}),\qquad \eta=(as,0^{,r-1}).
\]
For every odd $b$ and every odd positive integer $j$,
\begin{equation}\label{R31:odd_vanishing}
 m_{j\nu}\bigl(\mathcal T_{m,n,r}\bigr)=0
       \quad\hbox{in coordinate degree }sj.
\end{equation}
Equivalently the enhanced cone associated to this direction
has no invariant of odd degree. It contains no nonzero closed
orbit with the no-scaling property used in R-28--R-30.

In particular, for every odd $m\geq3$, take $s=r=2m+1$ and
$b=m$. The direction is strict, with $\nu_r>0$ and $\nu_1>as$,
but the same two-product common family cannot supply a stable
no-scaling point at this next row count. This is an obstruction
to that family and that all-integer method; it does not imply
vanishing or a multiplicity comparison for either full orbit
closure.

\begin{proof}
The $2m+1$ factors $L_1,\ldots,L_m,M_1,\ldots,M_m,\ell$
are linearly independent on a nonempty open set of the parameter
space. Their image is dense in $\mathcal T_{m,n,r}$. On this
open set a linear coordinate change identifies the form with
\[
 F=z^a(P+Q),\qquad P=x_1\cdots x_m,\quad Q=y_1\cdots y_m.
\]
The auxiliary extremal line for the one-row weight $\eta$ is
represented by $v^{as}$, where
\[
 v=\sum_{i=1}^m c_i x_i^*+\sum_{i=1}^m d_i y_i^*+e z^*.
\]
It is enough to consider the dense open set with $c_1c_2\ne0$.

Define an operator $h$ on the variable space by
\[
 h(x_1)=\frac{c_1}{c_2}x_2,\qquad
 h(x_2)=\frac{c_2}{c_1}x_1,
\]
and let it fix all other coordinates. Then $h^2=1$,
$\det h=-1$, $hP=P$, $hQ=Q$, and $hz=z$. Moreover
$v\circ h=v$, so the dual action also gives $hv=v$.
Thus $h$ fixes both $F$ and the auxiliary vector.

Choose $\zeta\in\mathbb C^*$ with $\zeta^r=-1$ and set
$g=\zeta h$. Then $g\in\mathrm{SL}_r$ and its action on
\[
 w=F^{\otimes s}\otimes v^{as}
\]
is
\[
 gw=\zeta^{ns-as}w=\zeta^{ms}w
                    =(\zeta^r)^b w=(-1)^b w.
\]
For odd $b$, every homogeneous invariant of odd degree must
therefore vanish at $w$. The generic factor and auxiliary
parameters considered above have dense image in the enhanced
cone, so every such invariant is identically zero on that
cone. The unconditional Cartan graded-coordinate extraction
in G-22 identifies its degree-$j$ invariant dimension with
$m_{j\nu}(\mathcal T_{m,n,r})$, proving
\eqref{R31:odd_vanishing}.

For the no-scaling consequence, let $w\ne0$ have closed orbit
in the enhanced cone. Every invariant takes the same value at
$w$ and $-w$, because its homogeneous odd-degree components
vanish identically. The cited closed-orbit separation theorem
(G-B15, Mukai2003 Theorem5.3 and Corollary5.4) then puts $w$
and $-w$ in the same $\mathrm{SL}_r$ orbit. Hence some group
element scales $w$ by $-1$, which contradicts the no-scaling
property.

Finally $\gcd(m,2m+1)=1$. The displayed choice $s=2m+1$,
$b=m$ satisfies the required equation, and $b$ is odd precisely
when $m$ is odd. The partition has size $as+rb=ns$, and its
first part exceeds $as$ by $b>0$. This proves the explicit
strict-direction obstruction claimed in the statement.
\end{proof}

\paragraph{Scope and distinction from the full cones.}
Restriction from either full coordinate ring to the common
subcone is a quotient. A zero multiplicity on the subcone
therefore gives no zero multiplicity on the full cone.
This result excludes the unchanged two-product no-scaling
strategy in the specified next-row directions, and leaves
other common pencils, even stretches, and other directions
open. No asymptotic comparison of the full permanent and
determinant multiplicities is asserted here.
